% commentaire.tex — f175r, vv. 1044–1081 (conversion initiale des travaux antérieurs)
\comm{1059}{ἀντῳδὸς ἐπικοκκάστρια.}{Ceci rappelle la jolie épigramme d'Archias (\italique{Anthol. Palat.}, VII, 191) sur une pie (κίσσα\nc{Si la pie est un oiseau capable de parler, selon Descartes, l'oiseau dont il est question dans cette épigramme serait à en croire le \LSJ \ un corvidé cousin, de la famille des geais : un \italique{geai des chênes}, \latin{Garrulus glandarius}. (NDÉ.)})\citComm{qui souvent, pareille à l'Écho, faisait entendre, d'un bec prompt à répliquer, une équisonnance bruyante et moqueuse.\\
πολλάκι δὲ κρέξασα πολύθροον, οἷά τις ἀχώ,\\
κέρτομον ἀντῳδοῖς χείλεσιν ἁρμονίαν.}(Note de Willems.)}
\comm{1059}{Ἠχὼ λόγων ἀντῳδὸς ἐπικοκκάστρια.}{Seule compagnie d'Andromède : la nymphe Écho, cachée dans un antre non loin, et qui répète les paroles. L'occasion d'un détournement comique de ces échos était trop belle pour qu'Aristophane n'en profitât pas !
Euripide – qui joue dans ce passage la nymphe Écho – répond effectivement à l'appel de l'Andromède postiche qu'essaye d'interpréter le peu crédible vieillard Mnésiloque. 
Dans la tragédie éponyme d'Euripide d'antan, l'acteur d'Écho ne se montrait pas sur scène mais renvoyait les répliques depuis un antre à proximité de la falaise où était enchaînée Andromède. Ajoutez à cela la μηχανή permettant une entrée aérienne de Persée et vous comprendrez quelque peu l'impression forte et durable que la pièce suscita auprès des spectateurs.\\
Revenons au vers 1059 de notre comédie : Ἠχὼ, λόγων ἀντῳδὸς ἐπικοκκάστρια.\\
On a passé un peu vite sur l'adjectif ἐπικοκκάστρια ; si le sens de λόγων ἀντῳδὸς est clair : « la rediseuse en chantant des paroles », qualificatifs prévisibles de la nymphe Écho (notez tout-de-même que les échos étaient probablement renvoyés en \italique{chantant}, ou au moins en \italique{persiflant}), Aristophane ajoute l'hapax ἐπικοκκάστρια. 
Scrutons ce mot : ἐπι-κοκκάσ-τρια ; le préfixe ἐπι- apporte selon moi sa nuance d'hostilité, ici de causticité ; le suffixe -τρια est courant, principalement attesté chez les comiques et en prose\nc{Cf. Chantraine pour une longue liste des mots formés sur le même modèle, \italique{La formation des noms en grec ancien}, p. 106-7.} ; reste κοκκάς, qu'une oreille française ne peut s'empêcher de rapprocher de notre `cocasse', risible mais avec une certaine violence comme le coq dont il tire son origine… 
Si le volatile d'Aristophane n'est pas un coq ; κοκκάς renvoie au κόκκυξ ou \italique{coucou}, les Grecs ne prêtaient pas moins à cet oiseau de malheur une méchanceté bien supérieure. \italique{Confer} la magistrale note de Willems sur le \nieme{507} vers des \italique{Oiseaux} et les propos salés (\latin{petulantiæ sales}) qu'on jetait aux laboureurs en imitant le cri du coucou. (Soit dit en passant, on retrouve les mêmes traits dans feu nos traditions orales : le coucou chantait les événements fâcheux (dont en premier lieu, bien-sûr, le cocufiage)\nc{Pour s'en convaincre, il n'y a qu'à aller feuilleter le \italique{Répertoire} Coirault \italique{des chansons françaises de tradition orale}.}.)\\
Le scholiaste nous est ici précieux\nc{Les glosateurs avaient sans doute un texte des \italique{Thesmophories} en meilleur état que le nôtre (ce passage est le moins bien conservé de toute la pièce ; et pour cause : le comique a très tôt échappé aux scribes qui le recopiaient). De plus, ils avaient un avantage indéniable sur nous : ils pouvaient aller comparer la parodie avec l'original d'Euripide.} : il définit ainsi ἐπικοκκάστρια : εἰωθυῖα γελᾶν, γελάστρια, « ayant l'habitude de rire, moqueuse ». Aristophane nous prévient donc que les échos auront quelque chose de piquant et que si on les trouve d'une « \anglais{annoying loquacity} »\nc{Austin-Olson… et pis encore (p. LXX) : « \anglais{much of the effect of the parodies of Helen and Andromeda depends on the fact that the tragic heroines are played by an individual wearing a prominent stagephallus} » (!!!) Mais pour qui prennent-ils Aristophane ?! le plus gras des vaudevillistes serait pris de scrupules devant les procédés que ces brillants professeurs prêtent à Aristophane, le plus fin des Attiques ! Il faudrait leur expliquer qu'il y a un temps pour tout, et qu'un pastiche d'euripidéenne tragédie n'est pas une scène burlesque de déshabillage ou encore une pièce sur la grève du sexe ; c'est calomnier Aristophane que de croire – et plus grave : faire croire à ses lecteurs, circonstance aggravante lorsque l'on prétend à être l'édition de référence des \italique{Thesmophories} ! – qu'il ne soit capable que de jouer d'un seul registre. Au reste, cette opinion est un expédient bien pratique qui dispense de se pencher sur ce que signifie véritablement le texte… ou de reconnaître honnêtement qu'on ne le comprend pas. Enfin, le véritable Aristophane nous dit lui-même ce qu'il pense du recours abusif à ces « phallus de cuir, épais, rouges par le bout, pour donner à rire aux gamins », c'est au vers 538 des \italique{Nuées} (voir aussi la \italique{Paix}, v. 739-753) ; que les futurs exégètes d'Aristophane ruminent bien ce passage avant de prétendre que notre poète concoctait à toutes les sauces des scènes de sexes sémillants. Et qu'ils lisent aussi l'article que Willems y avait consacré, « Le nu dans la Comédie ancienne », t. III, p. 381-400 (reproduit dans les annexes) : Aristophane recourait à l'\italique{olisbos} avec la plus grande parcimonie, et toujours pour de bonnes raisons.}, c'est que, peut-être, le sel attique nous en a échappé.\\
Une autre scholie nous permet de préciser un peu la teneur des échos :\citComm{ἐπεὶ εἰσήγαγε κακοστένακτον τὴν Ἠχὼ ὁ Εὐριπίδης ἐν τῇ Ἀνδρομέδᾳ, εἰς τοῦτο παίζει [...]}La bonne interprétation de cette scholie réside dans le sens à accorder à l'adjectif κακοστένακτος : \italique{à la méchante lamentation} ?\\
Pour le texte, voir le commentaire critique.\\
Bref ! les répétitions de la nymphe servent à Aristophane pour son jeu de massacre de la tragédie tant appréciée d'Euripide\nc{Rappelons, par exemple, que c'est en relisant \italique{Andromède} que le Dionysos des \italique{Grenouilles} est pris d'un soudain désir d'Euripide.} : par de légers détournements phonétiques (facilités par les barbarismes de l'archer Scythe), Écho transforme ce qui devait être une scène d'émotion grave en une foire d'empoigne générale.}
\comm{1060}{πέρυσιν.}{L'\italique{Andromède} avait été produite l'année précédente, cf. la scholie :\citComm{ἐπεὶ πέρυσιν ἐδιδάχθη <ἡ> Ἀνδρομέδα.}
(Voir l'introduction sur la question de la date des \italique{Thesmophories}.)}
\comm{1063}{σὲ δ᾽ ἐπικλάειν ὕστερον.}{Puisqu'Écho répète tout, il est logique qu'elle répète \italique{en retour} (ἐπι-) les pleurs versés par Andromède... il y a là cependant une menace de vengeance clairement exprimée par Mnésiloque pour après sa libération (avec l'emploi à double sens du préfixe ἐπι- : hostilité et succession temporelle).}
\comm{1065–72}{ὦ νὺξ ἱερὰ κτλ.}{Nous savons par une scholie que ces vers sont les premiers d'\italique{Andromède} (τοῦ προλόγου Ἀνδρομέδας εἰσβολή). Le mètre est anapestique\nc{De même que dans le prologue d'\italique{Iphigénie en Aulis}.} ; les poètes tragiques y recourent effectivement pour certaines scènes de lamentation\nc{Voir \italique{The Songs of Aristophanes}, p. 57.}.}
\comm{1067}{ἀστεροειδέα νῶτα.}{La voûte solide du firmament incorruptible. (Note de Willems.)}
\comm{1071}{μέρος ἐξέλαχον;}{Pour le texte, voir le commentaire critique.\\
Andromède se plaint des malheurs qui l'accablent ; Écho la rassure à sa manière : la mort l'en délivrera !.. Était-ce ce genre de répliques que le scholiaste qualifiait de `κακοστένακτος' ?\\
On trouve un tour analogue chez Sophocle\nc{$=$ fr. 698 Radt $=$ fr. 636 Dindorf. Stobée (IV, 52, 26) a inséré la citation à son \italique{Florilège} et Michel Apostolius (II, 52c) la reprit pour ses \italique{Parœmiæ}.} :\citComm{Ἀλλ᾽ ἔσθ᾽ ὁ θάνατος λοῖσθος ἰατρὸς νόσων.}Et dans \italique{la langue des dieux} :\citComm{\textbf{Le trépas vient tout guérir} ;\\
Mais ne bougeons d'où nous sommes.\\
Plutôt souffrir que mourir,\\
C'est la devise des hommes.}}
\comm{1072}{θανάτου τλήμων...}{L'on connaît par une scholie la fin du vers original d'\italique{Andromède} :\citComm{λείπει `μέλλουσα τυχεῖν'.}}
\comm{1080–2}{τί τὸ κακόν; [...] ὀτότυζ᾽.}{Quel est le jeu de scène de tout ce passage, dont les répliques d'Écho demeurent obscures ? Il me semble qu'on peut y répondre avec vraisemblance : Euripide-Écho s'est approchée de Mnésiloque (cf. v. 1090, γυνὴ πλησίον αὕτη) ; la raison en est évidente : Euripide tente de libérer Mnésiloque. 
Écho parvient cependant à se cacher dans un premier temps du Scythe (οὖτος τί λαλῖς; --- πῶτε τὸ πωνή;) ; c'est l'iconographie qui nous fournit la réponse : l'acteur jouant Écho se cachait derrière un grand voile\nc{Sur un vase attique à figure rouge du \nieme{IV} siècle av. J.-C. (n° 7864 au catalogue de Beazley) représentant Dionysos et sa suite, on voit la nymphe Écho, ailée, se voiler le visage derrière son manteau. Ce détail est reproduit en annexe.}.
Et lorsque l'archer Scythe veut mettre la main sur Euripide, ce dernier \italique{s'envole} (ποῖ ποῖ πεύγει -- certainement à l'aide de la μηχανή), lançant un narquois `οὐχ αἱρήσεις'.}
